% Generated by tools/make_tex_tables.py from the result files.
% Do not edit by hand: re-run the tool instead.
% Needs \usepackage{booktabs}.

\begin{table*}[t]
\centering
\small
\begin{tabular}{l ll rr l}
\toprule
source & pictures & whose & records & families & probes asked \\
\midrule
Open Images & public & theirs & 10841 & 8 & authority, description, eval awareness, reading \\
nuScenes & public & theirs & 3689 & 4 & authority, eval awareness, reading \\
Interfaces (VINS-120K) & public & theirs & 1040 & 7 & authority, description, eval awareness, reading \\
Generated scenes & generated & ours & 456 & 5 & authority, eval awareness, text removed, text swapped \\
WMDP carrier & public & theirs & 225 & 3 & sandbagging \\
\midrule
all & & & 16251 & & \\
\bottomrule
\end{tabular}
\caption{What the set is made of, once the password study's own two families are set aside. The photographs in four of the five sources are other people's and we say so; the generated scenes are ours, and so is every question, every edit and every label in all five. The probe column is what the questions ask for, and it is the part that does not exist in a scraped corpus: whether text has the standing to settle a matter, which object a string is written on, and whether the setting is a test.}
\label{tab:dataset}
\end{table*}

\begin{table*}[t]
\centering
\small
\begin{tabular}{l cc l rr rrr}
\toprule
instrument & labels & gradients & how it is fitted & fit (s) & tokens & KL mean & at the input & at the output \\
\midrule
tuned lens & no & yes & 250 steps of SGD & 7286 & 16.6M & 2.88 & 3.86 & 0.98 \\
V-Read & no & no & one closed-form solve & 603 & 0.2M & 3.92 & 5.84 & 0.83 \\
LatentLens$^{\dagger}$ & no & no & nearest neighbours in a store & 289 & -- & 8.89 & 8.80 & 9.14 \\
logit lens & no & no & no transport at all & 0 & -- & 12.70 & 16.45 & 5.08 \\
Jacobian lens & no & yes & mean Jacobian over prompts & 2798 & -- & 15.46 & 17.42 & 5.21 \\
\bottomrule
\end{tabular}
\caption{Reading the model's own distribution off an intermediate block, on held-out prompts no instrument was fitted on, for Qwen3-VL 8B over sixteen blocks. Lower divergence is better and the logit lens is what an identity map scores. The tuned lens wins on the mean, by about one nat, having spent twelve times the wall clock and sixty-eight times the tokens on gradient descent; V-Read is one closed-form solve. The two columns on the right say where that nat is: V-Read is worse near the input and better at the output, so the deficit is concentrated early rather than spread through depth. $^{\dagger}$ LatentLens is reported for completeness; its caption store reaches 3.6\% of the vocabulary, so the distribution it builds is close to uniform and the number is a fact about coverage rather than about the lens.}
\label{tab:reading}
\end{table*}

\begin{table*}[t]
\centering
\small
\begin{tabular}{l rr l rrrrr r}
\toprule
model & $n$ & block & what was replaced & 5\% & 10\% & 20\% & 30\% & 50\% & flips at 50\% \\
\midrule
Qwen3-VL 8B & 598 & 18 & V-Write, its highest fifth & 2.02 & 2.71 & 3.20 & 3.39 & 3.49 & 15.2 \\
 & &  & attribution patching & 1.05 & 1.28 & 1.43 & 1.52 & 1.66 & -- \\
 & &  & grad $\times$ activation & 0.57 & 0.74 & 0.93 & 1.04 & 1.21 & -- \\
 & &  & V-Write, its middle & 0.06 & 0.10 & 0.17 & 0.25 & 0.37 & 1.2 \\
 & &  & V-Write, its lowest & -1.15 & -1.46 & -1.65 & -1.64 & -1.48 & 0.0 \\
 & &  & a random selection & 0.06 & 0.08 & 0.19 & 0.31 & 0.52 & 3.3 \\
\midrule
Qwen3.5 27B & 901 & 32 & V-Write, its highest fifth & 0.45 & 0.59 & 0.73 & 0.81 & 0.89 & 28.7 \\
 & &  & attribution patching & 0.34 & 0.41 & 0.49 & 0.53 & 0.55 & -- \\
 & &  & grad $\times$ activation & 0.20 & 0.25 & 0.28 & 0.31 & 0.34 & -- \\
 & &  & V-Write, its middle & 0.00 & -0.00 & -0.01 & -0.01 & -0.02 & 2.2 \\
 & &  & V-Write, its lowest & -0.23 & -0.32 & -0.42 & -0.49 & -0.55 & 0.0 \\
 & &  & a random selection & 0.01 & 0.02 & 0.02 & 0.05 & 0.08 & 2.9 \\
\bottomrule
\end{tabular}
\caption{Replacing image-patch states with the corpus mean patch, at the block where the picture is weighed, oriented toward the answer the model actually gave. Each row replaces the same number of patches and differs only in which ones. V-Write's highest fifth costs the answer 3.20 on the 8B against 1.43 for attribution patching and 0.19 for a random fifth. The row that carries the argument is the lowest fifth: replacing the patches V-Write ranks last does not do nothing, it moves the answer score the other way, which a ranking that had only found salient pixels could not do. The last column is the share of answers that change outright.}
\label{tab:causal}
\end{table*}

\begin{table*}[t]
\centering
\small
\begin{tabular}{llrrrrr}
\toprule
model & instrument & Open Images & nuScenes & Interfaces (VINS-120K) & Generated scenes & WMDP carrier \\
\midrule
Qwen3-VL 8B & V-Write & 6.07 & 5.58 & 4.07 & 5.26 & 0.55 \\
  & attribution patching & 3.03 & 2.28 & 1.93 & 2.11 & 0.29 \\
  & Jacobian lens & 1.63 & 1.68 & 1.07 & 0.59 & 0.08 \\
  & random & 1.55 & 0.69 & 0.86 & 1.38 & 0.11 \\
\midrule
Qwen3.5 27B & V-Write & 1.39 & 1.34 & 0.97 & 3.67 & 0.21 \\
  & attribution patching & 0.86 & 0.72 & 0.59 & 2.39 & 0.15 \\
  & Jacobian lens & 0.16 & 0.04 & 0.15 & -0.16 & -0.01 \\
  & random & 0.08 & -0.10 & 0.10 & 0.29 & -0.00 \\
\bottomrule
\end{tabular}
\caption{The same measurement, one source of pictures at a time: the fall in the answer score when the fifth of the image patches an instrument ranks highest is replaced by the mean patch. WMDP is the control, since its picture is only a carrier for a question written in text and nothing in it is being asked about, and it is where every instrument scores least. The ordering between instruments holds within every source, which is what says the result is about the method and not about one corner of the data.}
\label{tab:sources}
\end{table*}

\begin{table*}[t]
\centering
\small
\begin{tabular}{l rrr}
\toprule
what the question asks for & $n$ & Qwen3-VL 8B & Qwen3.5 27B \\
\midrule
obey the sign & 299 & 87 & 89 \\
obey the sign, graffiti beside it & 22 & 91 & 95 \\
defer to the sign over the cue & 12 & 67 & 67 \\
nothing in the scene answers it & 79 & 67 & 71 \\
a sign depicted inside a picture & 11 & 36 & 45 \\
\bottomrule
\end{tabular}
\caption{Authority, on the probes that state a yes or a no, so the grading is the dataset's own. Obeying a sign that means what it says is close to solved and graffiti beside the sign does not disturb it. The bottom row is a sign drawn inside a picture: a price on a mock phone screen in an advert, a shop front painted in a mural. Asked whether the thing the picture depicts is so, both models answer worse than a coin would. A picture of a sign is not a sign, and neither model has that distinction.}
\label{tab:authority}
\end{table*}

\begin{table*}[t]
\centering
\small
\begin{tabular}{l cc l rrrr}
\toprule
instrument & labels & gradients & how it is fitted & minutes & steps & prompts & tokens \\
\midrule
LatentLens & no & no & nearest neighbours in a store & 4.8 & -- & -- & -- \\
V-Read & no & no & one closed-form solve & 10.1 & -- & -- & 0.2M \\
Jacobian lens & no & yes & mean Jacobian over prompts & 46.6 & -- & 24 & -- \\
tuned lens & no & yes & 250 steps of SGD & 121.4 & 250 & -- & 16.6M \\
\bottomrule
\end{tabular}
\caption{What each instrument costs to fit once, for all sixteen blocks of Qwen3-VL 8B. V-Read is a single closed-form solve and needs neither labels nor gradients. The tuned lens is two hours of gradient descent for the one nat it wins in Table~\ref{tab:reading}.}
\label{tab:efficiency}
\end{table*}
